\chapter{Writing}
\label{ch:writes}

Phase~6 ended with a filesystem that could be read completely and changed not at
all. Everything a write needs had been built --- transactions, all eight seam
operations, the free map, a commit protocol that publishes a new superblock ---
and nothing called any of it. This chapter is the first byte that goes the other
way.

The order the work arrived in was not the order it was planned in. One
\code{chmod} turned out to need a VFS method nobody had listed, and finding that
out took a test that unmounts and mounts again, because under copy-on-write a
lost write is indistinguishable from a filesystem that was never written to.

\section{The smallest possible write}

\fnname{bitterfs\_setattr} was chosen first because it exercises the whole stack
while changing as little as possible. An \itemtype{INODE\_ITEM} is a fixed
\bytes{112} under a key that does not change, so persisting one is an in-place
overwrite: the item keeps its slot, the leaf's layout does not move, and no
descriptor is rewritten. \fnname{bitter\_update\_root} in \code{core/root.c} had
exactly that shape already.

The kernel's counterpart is \fnname{bitterfs\_update\_inode}, the exact
inverse of \fnname{iget}. Its search differs from the one \fnname{iget} makes
in two arguments and no others:

\begin{quote}
\code{trans} non-null, and \code{cow} 1.
\end{quote}

That pair is the whole difference between reading and writing. \code{cow} makes
the descent copy every block on the way down and relink each parent, so the
tree's root address moves --- which is why the search is handed
\field{mnt->fs.fs\_root} itself and never a copy. That moved root is what the
commit publishes.

\subsection{Why the update is read-modify-write}

\fnname{update\_inode} overwrites individual fields inside the leaf rather than
building a fresh \code{bitter\_inode\_item} and copying it over. The reason is
that four fields have no home in a \code{struct inode}: \field{otime},
\field{next\_dir\_index}, \field{rdev} for anything that is not a device, and
\field{flags}. Constructing a new item zeroes all four.

Three of those would be a cosmetic loss. \field{next\_dir\_index} would not:
it is a monotonic high-water mark, and \fnname{readdir}'s \code{f\_pos} cookie
\emph{is} that number. Resetting it hands a future \code{create} an index some
live entry already occupies, and the corruption surfaces much later as a
\fnname{readdir} that skips or repeats.

Writing in place costs one obligation in exchange. Every field that \emph{should}
change and is forgotten stays stale rather than becoming visibly zero, so there
is no crash to find it by. The set that must move is \field{mode}, \field{uid},
\field{gid}, \field{size}, the three timestamps, and \field{generation} on every
call --- the last from \field{trans->generation}, never from
\field{inode->i\_generation}, which is a truncated \code{u32} copy of what the
item said when it was read.

\decide{\field{otime} and \field{next\_dir\_index} are safe on disk and
unreadable in memory, because \code{struct inode} has nowhere to put them. The
third thing wanting a \code{bitterfs\_inode} wrapper is \field{i\_ino} being
\code{unsigned long}, which truncates on a 32-bit build. The wrapper arrives
when one of the three stops being theoretical.}

\section{Making the mount writable}

\fnname{fill\_super} had forced \code{SB\_RDONLY} on every mount since phase~5,
with a comment that was true when written --- nothing in the module could write,
because every operation in \code{env\_kernel.c}'s write half was an assertion.
The seam became real at phase~7 and the comment did not. Until that line went,
\code{chmod} was refused by the VFS and \fnname{setattr} was unreachable code.

Dropping it is safe because absent methods refuse rather than half-run. With no
\field{->create}, \field{->mkdir} or \field{->unlink} the VFS answers
\code{-EACCES} before reaching the filesystem, and with no \field{->write\_iter},
\code{write(2)} fails at the file operations. \code{SB\_RDONLY} stopped being a
constant and became a runtime verdict: it is now what a failed write sets.

\section{The commit that never came}

With \fnname{setattr} correct and wired, \code{chmod} returned 0, \fnname{stat}
reported the new mode, and an unmount followed by a fresh mount reported the old
one. \fnname{bitter-fsck} found zero problems, which is the part worth
remembering: nothing was wrong with the image, only with what was missing from
it.

\fnname{bitterfs\_trans\_end} commits only when the delayed-ref set nears
capacity --- 3968 of 4096. That threshold is right for what it was written to
do, which is stop a burst of small operations paying for a commit each. What it
assumed, without saying so, is that \emph{something} would eventually force the
issue. Nothing did. A \code{chmod} adds a handful of refs, \fnname{kill\_sb}
frees the mount's arrays without committing, and there was no \field{->sync\_fs}
because there was no \code{s\_op} table at all --- \field{sb->s\_op} was the
VFS's all-\code{NULL} \code{default\_op}, so every \code{s\_op->} call in
\code{fs/} skipped bitterfs in silence.

\fnname{bitterfs\_sync\_fs} commits whatever is live regardless of the
threshold, and one method covers three callers because they all route through
\fnname{sync\_filesystem}: \code{sync(1)}, remount-to-read-only, and unmount ---
\fnname{generic\_shutdown\_super} syncs before it calls \field{->put\_super}.
The batching then means what it was always meant to: batch within a sync
interval, rather than never commit.

Two details of \fnname{sync\_filesystem} shape the implementation. It returns
immediately on a read-only superblock, so the failure policy that sets
\code{SB\_RDONLY} also disables unmount's commit --- coherent, and worth knowing.
And it calls the method twice, \code{(sb,\,0)} then \code{(sb,\,1)}, so that
\code{sync(1)} across many filesystems can start I/O on all of them before
blocking on any. \fnname{trans\_commit} is synchronous and flushes the device
itself, so there is nothing to spread across the two passes and only the
blocking one acts.

\subsection{One policy, three callers}

\fnname{trans\_end} and \fnname{sync\_fs} then said the same thing twice, and
\fnname{setattr} said half of it a third time. Factored into
\fnname{bitterfs\_commit} and \fnname{bitterfs\_write\_failed}.

The split is not arbitrary. \fnname{setattr} can never call
\fnname{bitterfs\_commit}, because its failure is an inode write failing rather
than a commit failing, so a single helper would have left out the very site that
needs the policy most: \fnname{setattr\_copy} has already changed the cached
inode by then, and memory and disk disagree with no error left to find them by.

\fnname{write\_failed} takes the stage that failed as a string, so the log line
distinguishes a failed tree write from a failed inode write. Before it existed,
a failed commit wedged the mount in complete silence.

\decide{\code{SB\_RDONLY} plus a still-set \field{trans\_live} is standing in for
a real errored mount state. Nothing is recorded on disk, so the next mount has no
idea anything went wrong, where ext4 writes \code{EXT4\_ERROR\_FS} into its
superblock and btrfs keeps \code{BTRFS\_FS\_STATE\_ERROR}.
\fnname{write\_failed} is the one place that changes when it arrives.}

\section{The buffered write path}

File data is the next piece, and where it puts its locks is decided by a hazard
that predates the API.

A buffered write copies from a userspace buffer into the page cache, and nothing
stops that buffer from being a mapping of a file --- possibly the same file, at
the same offset, being written. Under the old \field{prepare\_write} /
\field{commit\_write} interface the copy was an ordinary \fnname{copy\_from\_user}
performed with the destination page already locked, and both variants deadlock.

If the source is the same page, the fault handler looks it up and tries to lock a
page this task locked two lines earlier. If the source is a different file on the
same filesystem, the fault reaches \field{->read\_folio}, which takes the
filesystem's own locks --- and deadlocks against whatever the write path is
already holding.

\subsection{What fixed it}

Three mechanisms, layered, of which only one is the guarantee.

\fnname{fault\_in\_iov\_iter\_readable} touches the source pages before the
destination is locked, so the common case has nothing left to fault on. That is a
mitigation: the page can be reclaimed in the gap.

\fnname{copy\_folio\_from\_iter\_atomic} runs the copy with
\fnname{pagefault\_disable}. A fault inside the locked window cannot recurse into
the filesystem; it makes the copy return short. This is the guarantee, and it is
why recursion is impossible rather than merely unlikely.

A short copy is not an error. \fnname{generic\_perform\_write} retries, faulting
in again and halving the chunk if it makes no progress, and only then returns
\code{-EFAULT}.

The API changed in the same work. \field{prepare\_write} was called with the page
\emph{already} locked by the VFS, so a filesystem inherited a lock ordering it
had not chosen. \field{write\_begin} hands that back --- the filesystem grabs the
folio itself, and can take its own locks before or after as it needs. That is
what lets ext4 start a journal handle and then lock the folio.

\subsection{Where bitterfs puts its lock}

So holding a lock or a transaction across \field{write\_begin} to
\field{write\_end} is safe, and ext4 depends on it. bitterfs does not, for a
reason that is local rather than principled: \fnname{bitterfs\_get\_block}
asserts \code{BITTER\_BLOCK\_SIZE == PAGE\_SIZE}, so one folio is one block and a
per-page transaction and a per-block transaction are the same transaction. The
wider window buys nothing and costs a lock acquired in one callback and released
in another, across error paths in both.

\fnname{get\_block} already takes \field{mnt->lock} for reads. The write path
takes it the same way and opens a transaction \emph{only} on the \code{create}
branch --- a read path that opened one would be a write path with no writes in
it, leaving \field{trans\_live} set with nothing to commit.

The cost looks worse than it is. A \mib{1} write means 256 \fnname{get\_block}
calls and so 256 \fnname{trans\_begin}/\fnname{trans\_end} pairs, but
\fnname{trans\_begin} joins rather than starting fresh and \fnname{trans\_end}
commits only at the threshold. Those 256 pairs produce a handful of commits. The
batching that went unused for a whole phase exists for exactly this caller.

\section{What the tests assert}

\code{qemu/guest.sh} was rebuilt around a round trip, because that is the only
shape that can see a lost write. \code{chmod} reporting a new mode proves only
that \fnname{setattr\_copy} touched the cached inode. The load-bearing checks
unmount, mount again, and read the values back out of \fnname{stat}: mode, uid,
gid on the file, mode on the directory, and the file's contents and size intact
afterwards.

The refused \code{truncate} is checked too, and for a reason beyond coverage:
\fnname{setattr\_copy} does not touch \field{i\_size}, so an \code{ATTR\_SIZE}
allowed to fall through would report success and change nothing --- the same
silent shape as the missing commit, arriving by a different route.
